% basic_principles_axioms.tex -- Book I. Axioms (12 axioms, jemmybutton-matched).
% Text from jemmybutton byrne-en-latex.tex (CC-BY-SA).
% Build: lualatex basic_principles_axioms.tex  ->  basic_principles_axioms.pdf
\documentclass[12pt]{article}
\usepackage{geometry}
\geometry{a5paper, margin=1.5cm}
\usepackage[lining]{ebgaramond}
\usepackage{amssymb}
\usepackage{enumitem}
\usepackage{byrne}
\def\mpPre{u := 1cm; textLabels := false;}

% Helper: axiom block without figure (full-width centered)
\newcommand{\axiomtext}[2]{%
  \vspace{0.4cm}%
  \begin{center}%
    \textbf{#1.}\\[5pt]%
    #2%
  \end{center}%
}

\begin{document}

\begin{center}
  {\large\textbf{Book~I.}}\\[4pt]
  {\Large\textbf{Axioms}}
\end{center}

\vspace{0.4cm}
\noindent\rule{\linewidth}{0.4pt}
\vspace{0.2cm}

\axiomtext{I}{Magnitudes which are equal to the same are equal to each other.}

\axiomtext{II}{If equals be added to equals the sums will be equal.}

\axiomtext{III}{If equals be taken away from equals the remainders will be equal.}

\axiomtext{IV}{If equals be added to unequals the sums will be unequal.}

\axiomtext{V}{If equals be taken away from unequals the remainders will be unequal.}

\axiomtext{VI}{The doubles of the same or equal magnitudes are equal.}

\axiomtext{VII}{The halves of the same or equal magnitudes are equal.}

\axiomtext{VIII}{The magnitudes which coincide with one another, or exactly fill the same space, are equal.}

\axiomtext{IX}{The whole is greater than its part.}

\axiomtext{X}{Two straight lines cannot include a space.}

\axiomtext{XI}{All right angles are equal.}

% --- Axiom XII: parallel postulate (has figure) ---
\defineNewPicture{%
  pair A, B, C, D, E, F, G, H;%
  numeric s;%
  s := 3/2u;%
  A := (0, 0);%
  B := (4/3s, 0);%
  C := (0, s);%
  D := (4/3s, s);%
  E := (1/3s, 8/6s);%
  F := (xpart(E), -2/6s);%
  G = whatever[A, B] = whatever[E, F];%
  H = whatever[C, D] = whatever[E, F];%
  byAngleDefine(B, G, E, byred,    SOLID_SECTOR);%
  byAngleDefine(D, H, F, byyellow, SOLID_SECTOR);%
  draw byNamedAngleResized();%
  draw byLine(A, B, byblue,  SOLID_LINE, REGULAR_WIDTH);%
  draw byLine(C, D, byred,   SOLID_LINE, REGULAR_WIDTH);%
  draw byLine(E, F, byblack, SOLID_LINE, REGULAR_WIDTH);%
  draw byLabelLine(0)(AB, CD, EF);%
  draw byLabelsOnPolygon(E, H, D)(OMIT_FIRST_LABEL+OMIT_LAST_LABEL, 0);%
  draw byLabelsOnPolygon(B, G, F)(OMIT_FIRST_LABEL+OMIT_LAST_LABEL, 0);%
}

\vspace{0.5cm}
\noindent
\begin{minipage}[t]{0.34\linewidth}
  \vspace{0pt}\centering
  \drawCurrentPicture
\end{minipage}\hfill
\begin{minipage}[t]{0.62\linewidth}
  \vspace{0pt}\centering
  \textbf{XII.}\\[5pt]
  If two straight lines
  (\drawUnitLine{AB} and \drawUnitLine{CD})
  meet a third straight line (\drawUnitLine{EF})
  so as to make the two interior angles
  (\drawAngle{H} and \drawAngle{G})
  on the same side less than two straight angles,
  these two straight lines will meet if produced
  on that side on which the angles are less than
  two right angles.\\[8pt]
  \raggedright
  The twelfth axiom may be expressed in any of the
  following ways:
  \begin{enumerate}[leftmargin=*, topsep=2pt, itemsep=2pt]
    \item Two diverging straight lines cannot be both
          parallel to the same straight line.
    \item If a straight line intersect one of the two
          parallel straight lines it must also intersect
          the other.
    \item Only one straight line can be drawn through a
          given point, parallel to a given straight line.
  \end{enumerate}
\end{minipage}

\end{document}
